\documentclass[11pt]{article}
\usepackage[margin=1in]{geometry}
\usepackage{amsmath,amssymb,amsthm}
\usepackage[hidelinks]{hyperref}
\usepackage{graphicx}

\newtheorem{theorem}{Theorem}
\newtheorem{lemma}{Lemma}
\newtheorem{remark}{Remark}

\DeclareMathOperator{\Id}{Id}

\title{Nonquadratic Fixed Points of \(f=\tau f^*\)}
\author{FA Banach run: \texttt{fa\_banach\_001}}
\date{2026-06-28}

\begin{document}
\maketitle

\section*{Source Problem}

The source paper is:
\begin{quote}
A. N. Iusem, D. Reem, and S. Reich,
\emph{Fixed points of Legendre-Fenchel type transforms},
arXiv:1708.00464.
\end{quote}

Remark 15.1 asks whether the fixed point equation
\[
f(x)=\tau f^*(Ex+c)+\langle w,x\rangle+\beta
\]
always has a unique solution when \(E\) is positive definite. The authors
explicitly mention that even the simple case
\[
f(x)=\tau f^*(x),\qquad 0<\tau\neq 1,
\]
is unclear, and they suspect non-uniqueness.

\begin{center}
\includegraphics[width=0.92\textwidth]{figures/open_problem_crop.png}
\end{center}

\section*{Result}

\begin{theorem}
For every \(\tau>0\), \(\tau\neq1\), there exists a nonquadratic,
differentiable, proper, lower semicontinuous convex function
\[
f:\mathbb R\to\mathbb R
\]
such that
\[
f=\tau f^* .
\]
Thus the positive-definite uniqueness question in Remark 15.1 has a negative
answer already in the one-dimensional equation \(f(x)=\tau f^*(x)\).
\end{theorem}

\section*{Proof Intuition}

If \(f\) is differentiable and strictly convex with derivative \(h=f'\), then
\((f^*)' = h^{-1}\). Differentiating \(f=\tau f^*\) suggests the relation
\[
h(x)=\tau h^{-1}(x).
\]
Conversely, if an increasing homeomorphism \(h\) satisfies
\[
h^{-1}(x)=h(x)/\tau,
\]
then \(f(x)=\int_0^x h(s)\,ds\) satisfies \(f=\tau f^*\), after the harmless
normalization \(f(0)=0\).

The construction of \(h\) is easiest in logarithmic coordinates on
\((0,\infty)\). The equation becomes the problem of taking an increasing square
root of a translation. Translation has many non-linear increasing square roots:
conjugate the half-translation by a nontrivial periodic perturbation of the
identity.

\section*{Construction of the Derivative}

Let
\[
a=\log\tau \neq 0.
\]
Choose a real number \(\varepsilon\neq0\) so small that
\[
|\varepsilon|<\frac{|a|}{2\pi}.
\]
Define
\[
H(u)=u+\varepsilon\sin\left(\frac{2\pi u}{a}\right),\qquad u\in\mathbb R.
\]
Then \(H\) is a strictly increasing homeomorphism of \(\mathbb R\), and
\[
H(u+a)=H(u)+a.
\]
Set
\[
R=H^{-1}\circ T_{-a/2}\circ H,
\qquad T_b(u)=u+b.
\]
Since \(H\) commutes with \(T_a\), we have
\[
R^2=H^{-1}T_{-a}H=T_{-a}.
\]
Now define, for \(x>0\),
\[
h(x)=\exp(R(\log x)+a),
\]
set \(h(0)=0\), and extend \(h\) oddly:
\[
h(-x)=-h(x)\qquad (x>0).
\]

\begin{lemma}
The function \(h:\mathbb R\to\mathbb R\) is a strictly increasing homeomorphism
and satisfies
\[
h^{-1}(x)=\frac{h(x)}{\tau}\qquad(x\in\mathbb R).
\]
\end{lemma}

\begin{proof}
On \((0,\infty)\), the function \(x\mapsto \log x\), followed by \(R\), followed
by \(u\mapsto \exp(u+a)\), is a strictly increasing homeomorphism from
\((0,\infty)\) onto \((0,\infty)\). Since \(R(u)=u-a/2+O(1)\) as
\(u\to\pm\infty\), we have \(h(x)\to0\) as \(x\downarrow0\) and
\(h(x)\to\infty\) as \(x\to\infty\). The odd extension is therefore a strictly
increasing homeomorphism of \(\mathbb R\).

Let \(x=e^u>0\). Then
\[
\frac{h(x)}{\tau}
=\exp(R(u)+a-a)
=\exp(R(u)).
\]
Using \(R^2(u)=u-a\), we get
\[
h\left(\frac{h(x)}{\tau}\right)
=h(e^{R(u)})
=\exp(R(R(u))+a)
=\exp(u-a+a)
=x.
\]
Thus \(h(h(x)/\tau)=x\) for \(x>0\). By oddness the same identity holds for
all \(x\in\mathbb R\). Since \(h\) is bijective, this is equivalent to
\[
h^{-1}(x)=h(x)/\tau .
\]
\end{proof}

\section*{Proof of the Fixed Point Equation}

Define
\[
f(x)=\int_0^x h(s)\,ds,\qquad x\in\mathbb R.
\]
Since \(h\) is continuous and increasing, \(f\) is differentiable, convex, and
lower semicontinuous. Since \(h\) is onto, \(f\) is superlinear at both
infinities; in particular \(f^*\) is finite everywhere. Moreover \(f(0)=0\) and
\(0\) is the unique minimizer, so \(f^*(0)=0\).

Because \(f'=h\), the derivative of the convex conjugate is
\[
(f^*)'=h^{-1}.
\]
Therefore the function
\[
\Phi(x)=f(x)-\tau f^*(x)
\]
has derivative
\[
\Phi'(x)=h(x)-\tau h^{-1}(x)=0
\]
for every \(x\). Hence \(\Phi\) is constant. At \(x=0\),
\[
\Phi(0)=f(0)-\tau f^*(0)=0,
\]
so \(\Phi\equiv0\), which is exactly
\[
f=\tau f^*.
\]

Finally, \(f\) is not quadratic. If it were quadratic, then its derivative
\(h\) would be linear. On \((0,\infty)\), linearity of \(h\) would force
\(R(u)+a-u\) to be constant, equivalently \(R=T_{-a/2}\). But
\[
R=H^{-1}T_{-a/2}H
\]
equals \(T_{-a/2}\) only if \(H\) commutes with \(T_{-a/2}\). The chosen
\[
H(u)=u+\varepsilon\sin(2\pi u/a)
\]
does not commute with \(T_{-a/2}\) when \(\varepsilon\ne0\), since the sine term
changes sign under a half-period translation. Thus \(h\) is non-linear and
\(f\) is nonquadratic.

\section*{Verification Notes}

The proof gives a continuum of examples, one for each nontrivial periodic
perturbation \(H\). The construction is normalized by \(f(0)=0\); adding a
constant would destroy the equation unless the constant is zero.

The result is a counterexample to uniqueness in the positive-definite case
\[
E=\Id,\qquad c=w=0,\qquad \beta=0,\qquad 0<\tau\neq1.
\]
It does not attempt to classify all solutions of the general equation with
arbitrary \(E,c,w,\beta\).

\section*{Novelty Check}

On 2026-06-28, I searched for exact and close variants of the simple equation
and the source-paper question:
\[
\texttt{"f = tau f*" Legendre Fenchel nonunique solution},
\]
\[
\texttt{"f(x)=\textbackslash tau f\^*(x)" convex conjugate nonquadratic},
\]
\[
\texttt{"Fixed points of Legendre-Fenchel type transforms" "non-uniqueness" "tau"},
\]
and
\[
\texttt{"f=alpha f\^*" "convex conjugate"}.
\]
No explicit later answer to this positive-definite uniqueness question was
found in that bounded search.

\section*{References}

\begin{thebibliography}{9}
\bibitem{IRR}
A. N. Iusem, D. Reem, and S. Reich,
\emph{Fixed points of Legendre-Fenchel type transforms},
arXiv:1708.00464.
\end{thebibliography}

\end{document}
